% !TEX root = main.tex
\chapter{LABELLED STRUCTURES AND EXPONENTIAL GENERATING FUNCTIONS}
\label{chap:labelled structures and exponential generating functions}

\section{Labelled classes}

\section{Admissible labelled constructions}

\section{Surjections, set partitions, and words}

\section{Alignments, permutations, and related structures}

\section{Labelled trees, mappings, and graphs}
% Note II.19.
\noteblock{II.19}{Labelled hierarchies.}

We want to obtain
\begin{equation*}
    L(z) = T\left(\frac{1}{2} e^{z/2 - 1/2}\right) + \frac{z}{2} - \frac{1}{2}
\end{equation*}
from
\begin{equation*}
    L = z + e^{L} - 1 - L
\end{equation*}

\stephead{1}{Re-arranging the target equation.}
Collect the $L$-terms on one side:
\begin{equation*}
    2L + 1 - e^{L} = z
\end{equation*}
This still does not resemble $T(z) = ze^{T(z)}$, so we nudge it into Lambert-$W$ territory.

\stephead{2}{Utilizing the Lambert W identity.}
The Cayley tree function is really Lambert $W$ in disguise. In particular, $T(z) = ze^{T(z)}$ can be rewritten as
\begin{equation*}
    -T(z)e^{-T(z)} = -z \leftrightarrow z = -T(-z)e^{-T(-z)}
\end{equation*}
So $W(z) = -T(-z)$, equivalently $T(z) = -W(-z)$.

\stephead{3}{Substitution and solving.}
Return to the original equation and isolate the exponential:
\begin{equation*}
    e^{L} = 2L - z + 1
\end{equation*}
Whenever a variable appears both linearly and in the exponent, Lambert $W$ is usually nearby. So rewrite the equation in the form $f(L)e^{f(L)} = \text{something}$:
\begin{equation*}
    1 = (2L - z + 1)e^{-L}
\end{equation*}
To make the linear term match the exponent, set
\begin{equation*}
    u = \frac{z - 1}{2} - L.
\end{equation*}
Then
\begin{equation*}
    2u = z - 1 - 2L = -(2L - z + 1)
\end{equation*}
Substituting this into the previous equation,
\begin{equation*}
    1 = (-2u)e^{-L}
\end{equation*}
Also, since $u = \frac{z - 1}{2} - L$, we have $-L = u - \frac{z - 1}{2}$. Plugging that into the exponent gives:
\begin{equation*}
    1 = -2u \cdot e^{u - \frac{z - 1}{2}} = -2u \cdot e^{u} \cdot e^{-\frac{z - 1}{2}}
\end{equation*}
Now isolate $ue^{u}$:
\begin{equation*}
    ue^{u} = -\frac{1}{2} e^{\frac{z - 1}{2}}
\end{equation*}
Since $ue^{u} = v$ implies $u = W(v)$, and we know $W(v) = -T(-v)$:
\begin{equation*}
    u = W\left(-\frac{1}{2} e^{\frac{z - 1}{2}}\right) = -T\left(\frac{1}{2} e^{\frac{z - 1}{2}}\right)
\end{equation*}
Substitute $u$ back in:
\begin{equation*}
    \frac{z - 1}{2} - L = -T\left(\frac{1}{2} e^{\frac{z - 1}{2}}\right)
\end{equation*}
and solve for $L$:
\begin{equation*}
    L = T\left(\frac{1}{2} e^{\frac{z - 1}{2}}\right) + \frac{z - 1}{2}
\end{equation*}

% II.5.2.
\topicblock{II.5.2}{Mappings and associated graphs.}

Let $\mathcal{T}$ be the class of all Cayley trees. Then:
\begin{equation*}
    F = (1 - T)^{-1} \implies F_{n} = n^{n}
\end{equation*}
This follows neatly from Lagrange inversion. The standard form says
\begin{equation*}
    z = T / \phi(T) \implies [z^{n}] T(z) = \frac{1}{n} [w^{n - 1}] \phi(w)^{n}
\end{equation*}
More generally,
\begin{equation*}
    [z^{n}] H(T(z)) = \frac{1}{n} [w^{n - 1}] H'(w)\phi(w)^{n}
\end{equation*}
Taking $H(w) = (1 - w)^{-1}$ and $\phi(w) = e^{w}$ gives:
\begin{equation*}
    \begin{aligned}
        F_{n}
        &= n![z^{n}] F(z) \\
        &= n!\frac{1}{n} [w^{n - 1}] \left(\frac{1}{1 - w} \right)'(e^{w})^{n} \\
        &= (n - 1)![w^{n - 1}] \frac{e^{nw}}{(1 - w)^{2}} \\
        &= (n - 1)![w^{n - 1}] \sum_{k = 0}^{\infty}(k + 1)w^{k} \cdot \sum_{j = 0}^{\infty} \frac{n^{j}}{j!}w^{j} \\
        &= (n - 1)!\sum_{j = 0}^{n - 1}((n - 1 - j) + 1)\frac{n^{j}}{j!} \\
        &= (n - 1)!\sum_{j = 0}^{n - 1}(n - j)\frac{n^{j}}{j!} \\
        &= (n - 1)!\left(\sum_{j = 0}^{n - 1} n \frac{n^{j}}{j!} - \sum_{j = 0}^{n - 1} j \frac{n^{j}}{j!}\right) \\
        &= (n - 1)!\left(\sum_{j = 0}^{n - 1} \frac{n^{j + 1}}{j!} - \sum_{j = 1}^{n - 1} \frac{n^{j}}{(j - 1)!}\right) \\
        &= (n - 1)!\left(\sum_{j = 0}^{n - 1} \frac{n^{j + 1}}{j!} - \sum_{j = 0}^{n - 2} \frac{n^{j + 1}}{j!}\right) \\
        &= (n - 1)!\frac{n^{n}}{(n - 1)!} \\
        &= n^{n}
    \end{aligned}
\end{equation*}

Next consider $K = \log\left((1 - T)^{-1}\right)$. Then
\begin{equation*}
     K_{n} = n^{n - 1}Q(n)
\end{equation*}
where
\begin{equation*}
    Q_{n} = 1 + \frac{n - 1}{n} + \frac{(n - 1)(n - 2)}{n^{2}} + \cdots
\end{equation*}
From Note II.18 we already know
\begin{equation*}
    n![z^{n}] \frac{T(z)^{k}}{k!} = \binom{n - 1}{k - 1} n^{n - k} \implies
        n![z^{n}] \frac{T(z)^{k}}{k} = (k - 1)!\binom{n - 1}{k - 1} n^{n - k} = \frac{(n - 1)!}{(n - k)!} n^{n - k}
\end{equation*}
Therefore,
\begin{equation*}
    \begin{aligned}
        K_{n}
        &= n![z^{n}] K(z) \\
        &= n![z^{n}] \sum_{j = 1}^{\infty} \frac{T(z)^{j}}{j} \\
        &= \sum_{j = 1}^{\infty} n![z^{n}] \frac{T(z)^{j}}{j} \\
        &= \sum_{j = 1}^{\infty} \frac{(n - 1)!}{(n - j)!} n^{n - j} \\
        &= n^{n - 1} \sum_{j = 1}^{\infty} \frac{(n - 1)!}{(n - j)!n^{j - 1}} \\
        &= n^{n - 1}\left(1 + \frac{n - 1}{n} + \frac{(n - 1)(n - 2)}{n^{2}} + \cdots\right)
    \end{aligned}
\end{equation*}

\topicblock{II.5.2}{Constrained mappings.}

The class of maps without fixed points has exponential generating function:
\begin{equation*}
    \frac{e^{-T(z)}}{1 - T(z)}
\end{equation*}
and we want to show
\begin{equation*}
    n![z^{n}] \frac{e^{-T(z)}}{1 - T(z)} = (n - 1)^{n}
\end{equation*}
Again, Lagrange inversion does the bookkeeping. With $H(w) = \frac{e^{-w}}{1 - w}$ and $\phi(w) = e^{w}$, we obtain
\begin{equation*}
    \begin{aligned}
        n![z^{n}] \frac{e^{-T(z)}}{1 - T(z)}
        &= n!\frac{1}{n} [w^{n - 1}] \left(\frac{e^{-w}}{1 - w}\right)'(e^{w})^{n} \\
        &= (n - 1)![w^{n - 1}] \frac{we^{(n - 1)w}}{(1 - w)^{2}} \\
        &= (n - 1)![w^{n - 2}] \frac{e^{(n - 1)w}}{(1 - w)^{2}} \\
        &= (n - 1)![w^{n - 2}] \sum_{k = 0}^{\infty}(k + 1)w^{k} \cdot \sum_{j = 0}^{\infty} \frac{(n - 1)^{j}}{j!}w^{j} \\
        &= (n - 1)!\sum_{j = 0}^{n - 2}((n - 2 - j) + 1)\frac{(n - 1)^{j}}{j!} \\
        &= (n - 1)!\sum_{j = 0}^{m - 1}(m - j)\frac{m^{j}}{j!} \\
        &= (n - 1)!\frac{m^{m}}{(m - 1)!} \\
        &= (n - 1)!\frac{(n - 1)^{n - 1}}{(n - 2)!} \\
        &= (n - 1)^{n}
    \end{aligned}
\end{equation*}

% II.5.3.
\topicblock{II.5.3}{Unrooted trees and acyclic graphs.}

To see why $U(z) = \int_{0}^{z} \frac{T(w)}{w}\,\dd w$, start from the exponential generating functions
\begin{equation*}
    T(z) = \sum_{n \geq 0} T_{n} \frac{z^{n}}{n!}, \quad
    U(z) = \sum_{n \geq 0} U_{n} \frac{z^{n}}{n!}
\end{equation*}
Differentiate $U(z)$ and multiply by $z$; using $T_{0} = 0$, this gives:
\begin{equation*}
    z \frac{\dd}{\dd z} U(z) = z \frac{\dd}{\dd z} \sum_{n \geq 0} U_{n} \frac{z^{n}}{n!} = z \sum_{n \geq 1} nU_{n} \frac{z^{n - 1}}{n!}
        = \sum_{n \geq 1} T_{n} \frac{z^{n}}{n!} = T(z)
\end{equation*}
So $\frac{\dd}{\dd z} U(z) = \frac{T(z)}{z}$. Integrating from $0$ to $z$ gives:
\begin{equation*}
    U(z) - U(0) = \int_{0}^{z} \frac{T(w)}{w}\,\dd w
\end{equation*}
Now use the functional equation for rooted trees:
\begin{equation*}
    T(z) = ze^{T(z)}
\end{equation*}
This is equivalent to $z = Te^{-T}$. Differentiate with respect to $T$:
\begin{equation*}
    \dd z = \left(e^{-T} - Te^{-T}\right)\dd T = e^{-T}(1 - T)\dd T
\end{equation*}
Substituting $z = Te^{-T}$ and $\dd z$ into the integral turns it into
\begin{equation*}
    U(z) = \int \frac{T}{Te^{-T}} \cdot e^{-T}(1 - T)\,\dd T = \int (1 - T)\,\dd T = T - \frac{T^{2}}{2} + C
\end{equation*}
Since $U(0) = 0$ and $T(0) = 0$, the constant is $C = 0$.

% Note II.22.
\noteblock{II.22}{2–Regular graphs.}

The clean way to view this is as a bijection between clouds and 2-regular graphs on $n$ labelled vertices. Number the lines $1, 2, \dots, n$.
\begin{itemize}
    \item Line $i$ in the cloud $\leftrightarrow$ vertex $i$ in the graph.
    \item Intersection point $\{i,\; j\}$ in the cloud $\leftrightarrow$ edge $i - j$ in the graph.
\end{itemize}
Let $n_{i}$ be the number of chosen intersection points on line $i$.
Since each point lies on two lines, $\sum_{i = 1}^{n} n_{i} = 2n$.
\textit{If no three points are collinear}, then $n_{i} \leq 2$ for all $i$, so $\sum_{i = 1}^{n} n_{i} \leq 2n$.
Equality $\sum_{i = 1}^{n} n_{i} = 2n$ forces $n_{i} = 2$ for every $i$, so every vertex has degree 2.

Conversely, if every vertex has degree 2, then $n_{i} = 2$ for all $i$, so each line contains exactly two intersection points; in particular, no three chosen points are collinear.

So the two conditions are equivalent.

% Note II.23.
\noteblock{II.23}{Wright’s surgery.}

\stephead{1}{Deconstruct the graph.}

Start with a connected graph $G$ of excess $k + 1$. Choose an edge and orient it, so we have a source vertex $A$ and a target vertex $B$. Remove that directed edge.
At that point exactly two things can happen:
\begin{itemize}
    \item \textbf{Case 1 (Connected)}: The removed edge lay on a cycle, so the graph stays connected. We lost one edge and no vertices, hence the excess drops to $k$. What remains is a connected graph with two ordered distinguished vertices, $A$ and $B$.
    \item \textbf{Case 2 (Disconnected)}: The removed edge was a bridge, so the graph splits into two connected components $G_{1}$ and $G_{2}$. If $G_{1}$ has excess $h$, then $G_{2}$ has excess $k - h$. Each component carries one distinguished vertex, recording where the deleted edge used to attach.
\end{itemize}

\stephead{2}{Reconstruct the graph.}

To reconstruct all directed-edge-pointed graphs of excess $k + 1$, reverse the two cases:
\begin{itemize}
    \item \textbf{Reversing Case 1}: Start from one connected graph of excess $k$. Choose an ordered pair of distinct vertices $A,B$ and draw a directed edge from $A$ to $B$. If $A$ and $B$ were already adjacent, this would create a forbidden multi-edge, so those choices must be excluded.
    \item \textbf{Reversing Case 2}: Start from two independent connected components, of excess $h$ and $k - h$. Choose one vertex in each and join them by a directed edge from $A$ to $B$. This merges the two components into a connected graph of excess $k + 1$.
\end{itemize}

Now match these operations with the differential recurrence:
\begin{equation*}
    2\partial_{y}w_{k + 1} = \left(z^{2}\partial_{z}^{2}w_{k} - 2y\partial_{y}w_{k}\right) +
        \sum_{h = -1}^{k + 1} \left(z\partial_{z}w_{h}\right) \cdot \left(z\partial_{z}w_{k - h}\right)
\end{equation*}

\stephead{3}{Interpret the left-hand side.}
\begin{itemize}
    \item $w_{k + 1}$: Represents the universe of connected graphs of excess $k + 1$.
    \item $\partial_{y}$: Differentiation with respect to the edge-marker $y$ means ``choose an edge and delete it''.
    \item $2$: The deleted edge is oriented, so we must decide which endpoint is the source and which is the target.
    \item Together, $2\partial_{y}w_{k + 1}$ counts connected excess-$(k + 1)$ graphs with one remembered directed edge removed.
\end{itemize}

\stephead{4}{Interpret the connected term.}
\begin{itemize}
    \item $w_{k}$: The base graph of excess $k$.
    \item $\partial_{z}^{2}$: Differentiate twice with respect to the vertex-marker to select an ordered pair of distinct vertices $A,B$.
    \item $z^{2}$: Put the two marked vertices back into the count.
    \item $-2y\partial_{y}w_{k}$: This removes the bad choices where $A$ and $B$ were already joined by an edge. The factor 2 again accounts for orientation.
\end{itemize}

\stephead{5}{Interpret the disconnected term.}
\begin{itemize}
    \item $w_{h}$ and $w_{k - h}$: These are the two independent connected components, so their generating functions multiply.
    \item $z\partial_{z}$: This is the standard pointing operator for vertices.
        \begin{itemize}
            \item Applied to the first component, it selects the source vertex $A$.
            \item Applied to the second component, it selects the target vertex $B$.
        \end{itemize}
    \item $\sum_{h = -1}^{k + 1}$: Sum over all ways to distribute the remaining excess between the two components. The lower bound $-1$ appears because a tree already has excess $-1$.
\end{itemize}

\section{Additional constructions}
% Note II.28.
\noteblock{II.28}{Do two permutations generate the symmetric group?}

We want to show
\begin{equation*}
    n!^{2}\pi_{n} = (n - 1)!I_{n + 1}
\end{equation*}
The first useful observation is that the exponential generating function of pairs of permutations,
\begin{equation*}
    \sum_{n \geq 0} n!^{2} \frac{z^{n}}{n!} = \sum_{n \geq 0} n!z^{n},
\end{equation*}
is exactly the ordinary generating function $P(z)$ of permutations.
By the exponential formula, the exponential generating function $C(z)$ of connected graphs therefore satisfies
\begin{equation*}
    C(z) = \log P(z) = \log\left(\sum_{n \geq 0} n!z^{n}\right)
\end{equation*}
Since every permutation decomposes uniquely into a sequence of indecomposable blocks, the generating function $I(z)$ of indecomposable permutations satisfies
\begin{equation*}
    I(z) = 1 - \frac{1}{P(z)}
\end{equation*}
Now differentiate and extract the consequence:
\begin{equation*}
    C'(z) = \frac{P'(z)}{P(z)} = P'(z)(1 - I(z)) = P'(z) - P'(z)I(z)
\end{equation*}
Next we extract the coefficient of $z^{n - 1}$ from both sides.

\stephead{1}{Extract the left-hand side coefficient.}
Since $C(z) = \sum_{n \geq 1}(n!^{2}\pi_{n})\frac{z^{n}}{n!}$, its derivative is:
\begin{equation*}
    C'(z) = \sum_{n \geq 1}(n!^{2}\pi_{n})\frac{z^{n - 1}}{(n - 1)!}
\end{equation*}
So the coefficient of $z^{n - 1}$ is:
\begin{equation*}
    [z^{n - 1}] C'(z) = \frac{n!^{2}\pi_{n}}{(n - 1)!} = n \cdot n!\pi_{n}
\end{equation*}

\stephead{2}{Extract the right-hand side coefficient.}
Differentiate both sides of $P(z)(1 - I(z)) = 1$ with respect to $z$:
\begin{equation*}
    P'(z)(1 - I(z)) - P(z)I'(z) = 0 \implies P'(z) - P'(z)I(z) = P(z)I'(z)
\end{equation*}
So it remains to extract the coefficient of $z^{n - 1}$ from $P(z)I'(z)$:
\begin{equation*}
    [z^{n - 1}] (P'(z) - P'(z)I(z)) = [z^{n - 1}] P(z)I'(z) = \sum_{k = 1}^{n} kI_{k}(n - k)!
\end{equation*}
Call this sum $S_{n}$. The whole problem is now to show $S_{n} = I_{n + 1}$.
Extracting the coefficient of $z^{N}$ from $P(z) = 1 + I(z)P(z)$ for $N \geq 1$ gives the recurrence
\begin{equation} \label{eq:II.28.1}
    \sum_{k = 1}^{N} I_{k}(N - k)! = N!
\end{equation}
Apply \ref{eq:II.28.1} with $N = n + 1$:
\begin{equation} \label{eq:II.28.2}
    \sum_{k = 1}^{n + 1} I_{k}(n + 1 - k)! = (n + 1)! \implies \sum_{k = 1}^{n} I_{k}(n + 1 - k)! = (n + 1)! - I_{n + 1}
\end{equation}
Expand the left-hand side of \ref{eq:II.28.2}:
\begin{equation*}
    \begin{aligned}
        \sum_{k = 1}^{n} I_{k}(n + 1 - k)!
        &= \sum_{k = 1}^{n} I_{k}(n + 1 - k)(n - k)! \\
        &= (n + 1)\sum_{k = 1}^{n} I_{k}(n - k)! - \sum_{k = 1}^{n} kI_{k}(n - k)! \\
    \end{aligned}
\end{equation*}
By \ref{eq:II.28.1} with $N = n$:
\begin{equation*}
    \sum_{k = 1}^{n} I_{k}(n - k)! = n!
\end{equation*}
Hence the left-hand side of \ref{eq:II.28.2} is:
\begin{equation*}
    \sum_{k = 1}^{n} I_{k}(n + 1 - k)! = (n + 1)n! - S_{n} = (n + 1)! - S_{n}
\end{equation*}
\stephead{3}{Conclude the identity.}
Substituting back into \ref{eq:II.28.2} gives $S_{n} = I_{n + 1}$.

Finally, $n \cdot n!\pi_{n} = S_{n}$, so
\begin{equation*}
    n!^{2}\pi_{n} = (n - 1)!I_{n + 1}.
\end{equation*}

\section{Perspective}
